% ===== 5. SWOT: цифровий тренер =====
\begin{frame}{SWOT: цифровий тренер для автоспорту}
  \swotgrid
    {\item Не потрібне нове обладнання
     \item Відкритий формат GPMF і бібліотеки
     \item Демо: середня похибка кола 0,035 с}
    {\item GPS у телефоні: точність \textasciitilde{}4,9 м
     \item Криті картодроми без GPS
     \item Межі траси з відео поки не будуються}
    {\item Живий тренер — €250–350 на день
     \item Track Titan: 200\,000+ користувачів
     \item Ріст картингу: K1 Speed — 107 локацій}
    {\item Tenth Finder і Kart Track дають поради
     \item Дешевий RaceChrono Pro за \$19.99
     \item Камери без GPS, як HERO12}

  \note{[Доповідач 2] Тепер SWOT цифрового тренера. Головна сила — пілоту не треба купувати нове обладнання, досить камери чи телефона. GoPro зберігає телеметрію у відкритому форматі GPMF, і для нього вже є бібліотеки. А наше демо на реальній сесії показало, що час кола за GPS відрізняється від офіційного хронометражу в середньому на 0,035 секунди. Слабкі місця — точність GPS у телефоні, криті картодроми без сигналу і межі траси, які з відео поки не будуються. Можливості — дорогі живі тренери й інвестиції в Track Titan. Загрози — конкуренти, які вже дають поради, дешевий RaceChrono і те, що Камери без GPS, як HERO12.}
\end{frame}

% ===== 6. SWOT: розумний кошик =====
\begin{frame}{SWOT: розумний кошик для товарів}
  \swotgrid
    {\item Офіційний MCP: 40 інструментів
     \item Історія покупок онлайн і в магазинах
     \item Сам пропонує кошик, не чекає запиту}
    {\item «Що є вдома» — лише здогадка
     \item Кошик помиляється з фасуваннями
     \item Усі дані належать Сільпо}
    {\item 81\% купують продукти в магазинах
     \item Вибір страв: десятки хв на тиждень
     \item Відходи: 79 кг їжі на людину на рік}
    {\item Сільпо може зробити те саме
     \item Залежність від правил одного MCP
     \item Харчові дані можуть бути чутливими}

  \note{[Доповідач 2] SWOT розумного кошика. Сильна сторона — офіційний доступ до Сільпо: сорок інструментів, вхід через OAuth і приклади коду. До того ж видно історію покупок не лише онлайн, а й у магазинах. І продукт сам приходить з кошиком, а не чекає запиту. Але оцінка того, що є вдома, — лише здогадка, а в тестах автоматичний кошик уже помилявся з фасуваннями й наявністю. Власних даних у нас немає. Можливості — 81\% українців досі купують продукти в магазинах, тож онлайну є куди рости, плюс тема харчових відходів. Головні загрози — Сільпо може зробити те саме у своєму застосунку, а ми повністю залежимо від правил одного сервера.}
\end{frame}
